\section{The configuration file}
\label{sec:configuration}

The configuration file describes the hardware and the way to map onto it.
It has an optional \key{description} and five sections, each optional:
what a section leaves out keeps the default value of
\cref{tab:platform-keys,tab:exec-keys,tab:interference-keys,tab:network-keys,tab:method-keys}.
Without a \key{platform} section the platform is the built-in zonal one.

\begin{center}
\begin{tikzpicture}[node distance=2.5mm,
	sec/.style={box,minimum width=27mm,minimum height=13mm,
		font=\sffamily\scriptsize,text width=25mm}]
\node[sec] (p) {\textbf{platform}\\ECUs, roles, cores, memory};
\node[sec,right=of p] (e) {\textbf{execution\_model}\\scheduling, LET, cap};
\node[sec,right=of e] (i) {\textbf{interference\_model}\\shared memory};
\node[sec,right=of i] (n) {\textbf{network\_model}\\topology, TSN};
\node[sec,right=of n,fill=accentlight,draw=accent] (m) {\textbf{method}\\allocation, moves, sizing};
\node[above=2mm of $(p.north west)!0.5!(n.north east)$,font=\sffamily\scriptsize\bfseries,text=brand]
	{the system: what the hardware is and how it behaves};
\node[above=2mm of m.north,font=\sffamily\scriptsize\bfseries,text=accent!80!black]
	{how to map};
\end{tikzpicture}
\end{center}

\begin{tip}[Start from a template]
\cmd{mapping -{}-example-config T} prints a complete configuration file,
every key with its value, for the template \cmd{T}: \cmd{zonal} (the
built-in platform), \cmd{domain}, \cmd{hybrid} or \cmd{distributed}
(\cref{fig:templates}). \cmd{mapping -{}-config FILE -{}-print-config} prints
the configuration in effect, options included: the way to check what the
tool understood.
\end{tip}

\subsection{\texttt{platform}: the ECUs}
\label{sec:conf-platform}

\key{"platform": \{"ecus": [...]\}} lists 1 to 32 ECUs. The ECU order sets
the numbering of the cores (central computers first, then one core of each
ECU of a role at a time) and the order in which ECUs receive new cores.

\begin{shell}
\{"name": "HPC", "role": "central",\\
\ "cores": [\{"type": "big", "count": 4, "speed\_percent": 100\},\\
\ \ \ \ \ \ \ \ \ \ \{"type": "little", "count": 4, "speed\_percent": 40\}],\\
\ "added\_core\_type": "big", "memory\_kb": 524288, "max\_cores": 12\}
\end{shell}

\begin{table}[ht]
\centering\small\row
\caption{The keys of an ECU (\key{platform.ecus}).}
\label{tab:platform-keys}
\begin{tabularx}{\textwidth}{@{}>{\raggedright\arraybackslash}p{30mm}>{\raggedright\arraybackslash}p{26mm}X@{}}
\toprule\rowcolor{white}
key & default & meaning \\
\midrule
\key{name} & (mandatory) & at most 31 characters, distinct \\
\key{role} & (mandatory) & \key{central}, \key{domain}, \key{zone} or \key{ecu} (\cref{sec:concepts}) \\
\key{cores} & (mandatory) & one group per core type: \key{type} (\key{big}, \key{little}, \key{mcu}, \key{lockstep}), \key{count} (initial cores, may be 0) and \key{speed\_percent} (of the reference core). A type without a group does not exist on this ECU \\
\key{added\_core\_type} & the first type & type of the cores the tool adds to this ECU (it adds a core of another type of the ECU when a sub-chain may only run there) \\
\key{memory\_kb} & 0 (unlimited) & memory for the code and data of the tasks \\
\key{max\_cores} & 0 (unlimited) & largest number of cores: a fixed hardware variant \\
\bottomrule
\end{tabularx}
\end{table}

\begin{figure}[p]
\centering
\begin{tikzpicture}[node distance=3mm]
% ---- zonal
\node[font=\sffamily\small\bfseries,text=brand,anchor=west] (tz) at (0,0.55) {zonal};
\node[font=\sffamily\scriptsize,text=cmuted,anchor=west] at ([xshift=1mm]tz.east) {(built-in): a central computer, five zone controllers};
\node[switch] (zsw) at (6.4,-1.6) {SW};
\node[ecucentral,minimum width=21mm] (zpcu) at (6.4,-0.45) {PCU\cores{8 big 100\,\%}};
\foreach \i/\x in {1/0.9,2/3.65,3/6.4,4/9.15,5/11.9} {
	\node[ecuzone,minimum width=21mm] (zz\i) at (\x,-2.9) {ZCU\i\cores{4 MCU 50\,\%}};
	\draw[link] (zz\i.north) -- (zsw);
}
\draw[link] (zpcu) -- (zsw);
% ---- domain
\node[font=\sffamily\small\bfseries,text=teal,anchor=west] (td) at (0,-4.5) {domain};
\node[font=\sffamily\scriptsize,text=cmuted,anchor=west] at ([xshift=1mm]td.east) {five domain controllers, no central computer};
\node[switch] (dsw) at (6.4,-6.1) {SW};
\foreach \n/\c/\x/\y [count=\k] in {ADAS/4 big 100\,\%/0.9/-5.2,Infotainment/2 big + 4 little/3.65/-5.2,
	Powertrain/3 MCU + 1 lock-step/9.15/-5.2,Chassis/2 MCU + 2 lock-step/11.9/-5.2} {
	\node[ecudomain,minimum width=24mm] (dd\k) at (\x,\y) {\n\cores{\c}};
	\draw[link] (dd\k.south) -- (dsw);
}
\node[ecudomain,minimum width=24mm] (dBody) at (6.4,-7.25) {Body\cores{2 MCU 50\,\%}};
\draw[link] (dBody) -- (dsw);
% ---- hybrid
\node[font=\sffamily\small\bfseries,text=brand,anchor=west] (th) at (0,-8.15) {hybrid};
\node[font=\sffamily\scriptsize,text=cmuted,anchor=west] at ([xshift=1mm]th.east) {a central computer, an ADAS domain controller, four zones};
\node[switch] (hsw) at (6.4,-10.3) {SW};
\node[ecucentral,minimum width=21mm] (hpcu) at (5.0,-9.15) {PCU\cores{8 big 100\,\%}};
\node[ecudomain,minimum width=21mm] (hadas) at (7.8,-9.15) {ADAS\cores{4 big 100\,\%}};
\draw[link] (hpcu.south) -- (hsw); \draw[link] (hadas.south) -- (hsw);
\foreach \n/\x [count=\k] in {ZCU\_FL/2.25,ZCU\_FR/5.0,ZCU\_RL/7.8,ZCU\_RR/10.55} {
	\node[ecuzone,minimum width=21mm] (hz\k) at (\x,-11.45) {\n\cores{4 MCU 50\,\%}};
	\draw[link] (hz\k.north) -- (hsw);
}
% ---- distributed
\node[font=\sffamily\small\bfseries,text=cmuted,anchor=west] (tdi) at (0,-13.1) {distributed};
\node[font=\sffamily\scriptsize,text=cmuted,anchor=west] at ([xshift=1mm]tdi.east) {twelve ECUs, one function each, at most two cores (four for ADAS), 100 Mbit/s};
\node[switch] (disw) at (6.4,-15.2) {SW};
\foreach \n/\x/\y in {ADAS/0.9/-14.0,ESC/2.75/-14.0,EPS/4.6/-14.0,ECM/8.2/-14.0,TCM/10.05/-14.0,BCM/11.9/-14.0,
	HVAC/0.9/-16.4,IPC/2.75/-16.4,DOOR\_FL/4.6/-16.4,DOOR\_FR/8.2/-16.4,DOOR\_RL/10.05/-16.4,DOOR\_RR/11.9/-16.4} {
	\node[ecuplain,minimum width=16mm,minimum height=7mm] (di) at (\x,\y) {\n};
	\draw[link] (di) -- (disw);
}
\end{tikzpicture}
\caption{The four templates of \opt{example-config}, star topology. The
zone, domain and distributed ECUs are wired to the sensors and actuators of
generated applications; ESC and EPS have one lock-step core each, ADAS two
big cores, the others one microcontroller core at 50\,\%.}
\label{fig:templates}
\end{figure}

\subsection{\texttt{execution\_model}: scheduling and communication}
\label{sec:conf-exec}

\begin{table}[ht]
\centering\small\row
\caption{The keys of \key{execution\_model}.}
\label{tab:exec-keys}
\begin{tabularx}{\textwidth}{@{}>{\raggedright\arraybackslash}p{36mm}>{\raggedright\arraybackslash}p{30mm}X@{}}
\toprule\rowcolor{white}
key & default & meaning \\
\midrule
\key{scheduling} & \key{fixed-priority preemptive} & the only one supported; written for the reader of the file \\
\key{communication} & \key{implicit} & between cores: \key{implicit} (an output is copied at the first LET instant after its producer ends) or \key{LET} (strict LET: published at the end of its producer's period) (\cref{sec:concepts-let}) \\
\key{utilization\_cap} & 0.8 & $\varphi$: largest load of a core, LET task included \\
\key{let\_task\_wcet\_us} & 50 & WCET of the LET task of a core that communicates, on the reference core \\
\key{wcet\_margin\_percent} & 0 & every WCET is raised by this percentage before the mapping \\
\bottomrule
\end{tabularx}
\end{table}

\subsection{\texttt{interference\_model}: shared-memory contention}
\label{sec:conf-interference}

The cores of one ECU share its memory; a request waits for the requests of
the other cores, served in round robin. A task then runs for its WCET plus
its requests times $(n-1)$ times the latency, $n$ the cores of the ECU that
run tasks; the tool iterates until the number of competing cores is the
one of the mapping.

\begin{table}[ht]
\centering\small\row
\caption{The keys of \key{interference\_model}.}
\label{tab:interference-keys}
\begin{tabularx}{\textwidth}{@{}>{\raggedright\arraybackslash}p{46mm}>{\raggedright\arraybackslash}p{24mm}X@{}}
\toprule\rowcolor{white}
key & default & meaning \\
\midrule
\key{shared\_memory\_arbitration} & \key{round-robin} & the only one supported \\
\key{memory\_latency\_ns} & 0 & time to serve one request; 0: no contention \\
\key{memory\_accesses\_per\_ms} & 1000 & requests per ms of reference WCET, for the tasks without \key{memory\_accesses} and for the LET tasks \\
\bottomrule
\end{tabularx}
\end{table}

\subsection{\texttt{network\_model}: the TSN backbone}
\label{sec:conf-network}

Every ECU has one full-duplex link to a switch. The \key{topology} gives
the switches: \key{star} (one switch), \key{tree} (one switch per ECU, all
linked to the switch of the root ECU), \key{ring} (one switch per ECU, in
ECU order) or \key{custom} (\cref{fig:topologies}). Routes are shortest
paths; the links of the central computers run at the backbone rate.

\begin{figure}[ht]
\centering
\begin{tikzpicture}[sw/.style={switch,minimum width=5mm,minimum height=4mm,font=\sffamily\tiny\bfseries},
	e/.style={draw=cmuted,fill=white,circle,inner sep=1.2pt,minimum size=3.2mm,font=\sffamily\tiny}]
% star
\begin{scope}
\node[font=\sffamily\scriptsize\bfseries,text=brand] at (0,1.7) {star};
\node[sw] (s) at (0,0) {};
\foreach \a/\n in {90/C,18/1,-54/2,-126/3,162/4} {
	\node[e] (e\n) at (\a:1.15) {\n}; \draw[link] (s) -- (e\n);}
\end{scope}
% tree
\begin{scope}[xshift=3.6cm]
\node[font=\sffamily\scriptsize\bfseries,text=brand] at (0,1.7) {tree};
\node[sw] (r) at (0,0.35) {};
\node[e] (rc) at (0,1.15) {C}; \draw[link] (r) -- (rc);
\foreach \x/\n in {-1.2/1,-0.4/2,0.4/3,1.2/4} {
	\node[sw] (t\n) at (\x,-0.55) {}; \draw[fast] (r) -- (t\n);
	\node[e] (te\n) at (\x,-1.3) {\n}; \draw[link] (t\n) -- (te\n);}
\end{scope}
% ring
\begin{scope}[xshift=7.2cm]
\node[font=\sffamily\scriptsize\bfseries,text=brand] at (0,1.7) {ring};
\foreach \a/\n in {90/C,18/1,-54/2,-126/3,162/4} {
	\node[sw] (q\n) at (\a:0.75) {};
	\node[e] (qe\n) at (\a:1.3) {\n}; \draw[link] (q\n) -- (qe\n);}
\draw[fast] (qC) -- (q1) -- (q2) -- (q3) -- (q4) -- (qC);
\end{scope}
% custom
\begin{scope}[xshift=10.9cm]
\node[font=\sffamily\scriptsize\bfseries,text=brand] at (0,1.7) {custom};
\node[sw] (u1) at (-0.7,0.5) {}; \node[sw] (u2) at (0.7,0.5) {};
\node[sw] (u3) at (0,-0.55) {};
\draw[fast] (u1) -- (u2) -- (u3) -- (u1);
\node[e] (ue1) at (-1.35,1.1) {C}; \draw[link] (u1) -- (ue1);
\node[e] (ue2) at (1.35,1.1) {1}; \draw[link] (u2) -- (ue2);
\node[e] (ue3) at (1.35,0.0) {2}; \draw[link] (u2) -- (ue3);
\node[e] (ue4) at (-0.5,-1.3) {3}; \draw[link] (u3) -- (ue4);
\node[e] (ue5) at (0.5,-1.3) {4}; \draw[link] (u3) -- (ue5);
\end{scope}
\end{tikzpicture}
\caption{The topologies, for a central computer C and four other ECUs:
the switches are dark, the backbone links thick. A custom topology may
have switches without ECUs and links of their own rates.}
\label{fig:topologies}
\end{figure}

A custom topology names its switches with the ECUs on each one, and the
links between switches (full duplex, at the backbone rate unless they
give theirs):
\begin{shell}
"topology": "custom",\\
"switches": [\{"name": "SW\_CENTRAL", "ecus": ["HPC"]\},\\
\ \ \ \ \ \ \ \ \ \ \ \ \{"name": "SW\_FRONT", "ecus": ["ADAS\_DC", "ZC\_FL", "ZC\_FR"]\},\\
\ \ \ \ \ \ \ \ \ \ \ \ \{"name": "SW\_REAR", "ecus": ["ZC\_RL", "ZC\_RR"]\}],\\
"switch\_links": [\{"from": "SW\_CENTRAL", "to": "SW\_FRONT", "mbps": 10000\},\\
\ \ \ \ \ \ \ \ \ \ \ \ \ \ \ \ \{"from": "SW\_CENTRAL", "to": "SW\_REAR", "mbps": 2500\},\\
\ \ \ \ \ \ \ \ \ \ \ \ \ \ \ \ \{"from": "SW\_FRONT", "to": "SW\_REAR"\}]
\end{shell}
The tool checks that every ECU is on one switch, that every switch can
reach every other one, and that no route is longer than 16 ports. With
FRER (ring or custom), the second route of a replicated message is the
shortest one that shares no link and no switch with the first between the
switch of the source and the switch of the destination.

\begin{table}[ht]
\centering\small\row
\caption{The keys of \key{network\_model}.}
\label{tab:network-keys}
\begin{tabularx}{\textwidth}{@{}>{\raggedright\arraybackslash}p{44mm}>{\raggedright\arraybackslash}p{24mm}X@{}}
\toprule\rowcolor{white}
key & default & meaning \\
\midrule
\key{standard} & \key{TSN} & the only one supported \\
\key{topology} & \key{star} & \key{star}, \key{tree}, \key{ring} or \key{custom} \\
\key{tree\_root} & first central & tree: the ECU whose switch is the root \\
\key{switches}, \key{switch\_links} & & custom topology only (above) \\
\key{link\_mbps} & 1000 & rate of the ECU links, Mbit/s \\
\key{backbone\_mbps} & link rate & rate of the links between switches and of the links of the central computers \\
\key{ecu\_link\_mbps} & & rates of single ECU links, e.g.\ \key{\{"HPC": 10000\}} \\
\key{switch\_latency\_us} & 2 & fabric and propagation delay per switch crossed \\
\key{link\_limit\_percent} & 100 & design limit of the LET load of a port; a port above it is overloaded \\
\key{traffic\_shaping} & \key{CBS} & \key{CBS}, \key{TAS} or \key{ATS} \\
\key{traffic\_classes} & 1 & 1 to 3 LET classes, by the period of their producer \\
\key{class\_max\_periods\_us} & 5000, 20000 & largest producer period of each class but the last (classes $-1$ values) \\
\key{idle\_slopes\_percent} & 75 (one class) & CBS: share of the port rate of each class \\
\key{proportional\_idle\_slopes} & true & CBS, several classes: the total of the shares divided at each port in proportion to the traffic of each class \\
\key{tas\_cycle\_us} & 500 & TAS: gate cycle \\
\key{tas\_windows\_percent} & 70 (one class) & TAS: window of each class, \% of the cycle \\
\key{frame\_preemption} & false & preemption of the best-effort frames \\
\key{frer} & false & ring or custom: messages of replicated chains over two disjoint routes \\
\bottomrule
\end{tabularx}
\end{table}

\subsection{\texttt{method}: how to map}
\label{sec:conf-method}

\begin{table}[ht]
\centering\small\row
\caption{The keys of \key{method}.}
\label{tab:method-keys}
\begin{tabularx}{\textwidth}{@{}>{\raggedright\arraybackslash}p{34mm}>{\raggedright\arraybackslash}p{18mm}X@{}}
\toprule\rowcolor{white}
key & default & meaning \\
\midrule
\key{allocation} & \key{FFDU} & BFDU (best fit: fullest core), WFDU (worst fit: emptiest core), FFDU (first fit), CADU (fewest network hops, then a central computer, then best fit); all by decreasing utilization \\
\key{sigma} & 0.4 & $\sigma$: largest utilization of a sub-chain, fraction of a reference core (at most the cap) \\
\key{affinity} & true & a sub-chain first tries the cores of the other sub-chains of its chain \\
\key{consolidation} & true & after the allocation, empty the cores whose sub-chains fit elsewhere (not with WFDU) \\
\key{swap\_moves} & false & also exchange sub-chains of two unschedulable cores, before the relocation moves (in the experiments it never repaired a core) \\
\key{relocation\_moves} & true & move sub-chains of unschedulable cores to schedulable ones \\
\key{local\_search} & true & move sub-chains to bring late task-chains on time \\
\key{replica\_separation} & \key{ecu} & replicas on different \key{ecu}s, or only on different \key{core}s \\
\key{fixed\_hardware} & false & add no core: a sub-chain that fits nowhere makes the mapping infeasible (software updates of a frozen platform) \\
\key{target\_scaling} & 0 (none) & sizing: lower the cap until the critical scaling factor reaches this value (e.g.\ 1.3) \\
\bottomrule
\end{tabularx}
\end{table}

\begin{warn}[Which rule?]
FFDU and BFDU give the fewest cores; CADU the fewest network messages, at
the price of more cores; WFDU spreads the load on every core, which is
wasteful at design time but good on hardware that is already paid for
(use case~4, \cref{sec:uc4}). \cmd{-{}-algorithm all} compares the four.
\end{warn}
